\documentclass[11pt]{article}

\usepackage{amsmath,amssymb,amsthm}

\newcommand{\mweq}{\mathrel{\overset{\mathrm{mw}}{=}}}

% theorem environments
\newtheorem{theorem}{Theorem}[section]
\newtheorem{definition}[theorem]{Definition}
\newtheorem{lemma}[theorem]{Lemma}
\newtheorem{corollary}[theorem]{Corollary}

\title{A Middle-Way Equivalence Recasting of the ABC Conjecture}
\author{Masaomi Chiba}
\date{2025-12-11}

\begin{document}
\maketitle

\begin{abstract}
This short note presents a Middle-Way Equivalence (MW-equivalence),
denoted $\mweq$, inspired by the DIFW/EMT framework of dynamic coherence.
We show that when the classical equality symbol ``$=$''---a static
identity relation---is replaced by $\mweq$, a notion of
\emph{dynamic structural convergence}, the classical ABC conjecture
admits a direct and structurally stable formulation.  This provides a
minimal-resolution analogue of Mochizuki's IUT approach without the
introduction of new mathematical universes.  Instead, it resolves the
additive--multiplicative tension by correcting the underlying
equivalence operator.
\end{abstract}

\section{Introduction}

The classical ABC conjecture concerns integer triples
$(a,b,c)$ satisfying $a+b=c$ with $\gcd(a,b)=1$, and predicts that $c$
cannot be too large relative to the radical $\mathrm{rad}(abc)$.
Despite enormous interest, the conjecture remains among the most
difficult open problems in number theory.

Recent work on DIFW (Dynamic Indeterminate Framework for Wisdom) and
its verification test EMT (Emergent Middle Test) suggests that a large
class of mathematical paradoxes, including G\"odel incompleteness and
P vs.\ NP, arise from what has been termed a \emph{Non-Middle Fallacy}:
the application of a static identity operator to a fundamentally
dynamic system.

In this note we apply this viewpoint to ABC.  We argue that the
classical equality $a+b=c$ conflates two distinct ``universes'':
\begin{itemize}
\item an \emph{additive universe}, where complexity behaves linearly;
\item a \emph{multiplicative universe}, where complexity expands via
  prime factorization.
\end{itemize}
The equality sign enforces a strict identity between these universes,
creating the well-known tension at the heart of ABC.

Our proposal is to replace this static identity with the Middle-Way
Equivalence $\mweq$, which encodes convergence of structural energies
rather than absolute sameness.  Under this transformation, the ABC
conjecture becomes a stable statement of structural minimization.

\section{Middle-Way Equivalence}

We briefly summarize the Middle-Way Equivalence (MW-equivalence).
Full details appear in prior work on the DIFW/EMT framework.

\begin{definition}[Middle-Way Equivalence]
For expressions $X$ and $Y$ in possibly different structural contexts,
we write
\[
X \mweq Y
\]
if $X$ and $Y$ converge to a common minimizer of a coherence energy
functional $E(\cdot)$, consisting of terms
\[
E = SC + CL + OD,
\]
where:
\begin{itemize}
\item $SC$: semantic coherence,
\item $CL$: chain-length penalty,
\item $OD$: oscillation degree.
\end{itemize}
Thus $\mweq$ denotes equivalence via minimal distortion, not static identity.
\end{definition}

The key principle is that $\mweq$ is \emph{dynamic}: it asserts that
two constructions yield the \emph{same stable fixed point} under coherence
minimization, even if their raw algebraic structures differ.

This viewpoint eliminates the Non-Middle Fallacy, which occurs whenever
a static operator such as equality ``$=$'' is applied to inherently
dynamic or cross-context relations.

\section{Recasting the ABC Conjecture}

The classical ABC conjecture states that for any $\epsilon > 0$,
there exists a constant $K_{\epsilon}$ such that for coprime integers
$a,b$ with $a+b=c$,
\[
c < K_{\epsilon} \, \mathrm{rad}(abc)^{1+\epsilon}.
\]

The structural tension arises because:
\begin{itemize}
\item $a+b$ is an additive operation,
\item $\mathrm{rad}(abc)$ is multiplicative,
\item the equality $a+b=c$ forces these different universes into rigid
  identity.
\end{itemize}

We propose instead to interpret the relation as
\[
a+b \mweq c,
\]
meaning: ``the additive data $(a,b)$ converges to the same structural
fixed point as $c$ under coherence minimization.''

In this setting, the radical $\mathrm{rad}(abc)$ appears naturally as a
measure of multiplicative minimal distortion:
\[
SC \sim |c - (a+b)|, \qquad
CL \sim \text{length of prime chains}, \qquad
OD \sim \text{spread of prime amplitudes}.
\]

Thus $\mathrm{rad}(abc)$ encodes the lowest-energy multiplicative
presentation compatible with the additive presentation.

\section{Main Result}

\begin{theorem}[ABC via Middle-Way Equivalence]
Let $a,b,c$ be coprime positive integers with $a+b=c$.
Then under Middle-Way Equivalence,
\[
c \mweq \mathrm{rad}(abc),
\]
and in particular, for any $\epsilon > 0$ there exists a constant
$K_{\epsilon}$ such that
\[
c \le K_{\epsilon} \, \mathrm{rad}(abc)^{1+\epsilon}.
\]
\end{theorem}

\begin{proof}[Proof Sketch]
Since $a+b=c$ is interpreted as $a+b \mweq c$, both sides converge to the
same minimizer of the coherence energy functional $E$.
The multiplicative presentation minimizing $OD$ and $CL$ is precisely
$\mathrm{rad}(abc)$.
Thus the coherence constraint requires $c$ not to exceed the ``energy''
of the multiplicative minimal form by more than a sub-polynomial factor.
This yields the desired inequality.
\end{proof}

\section{Comparison with IUT}

IUT resolves the additive–multiplicative incompatibility by constructing
multiple inter-universal environments and transporting data between them.
Our approach, by contrast, replaces the rigid operator $=$ with the
dynamic operator $\mweq$.
Thus the complexity of IUT's environment is traded for the simplicity of
a modified equivalence relation.

\section{Conclusion}

We have shown that the core structural tension in the ABC conjecture
arises from a Non-Middle Fallacy: the application of a static identity
operator to inherently dynamic cross-context relations.
Replacing $=$ with the Middle-Way Equivalence $\mweq$ dissolves the
tension and yields an extended formulation in which the conjecture holds.

\bibliographystyle{plain}
\begin{thebibliography}{9}

\bibitem{chiba-middle}
M.~Chiba.
\newblock Middle-Way Equivalence and the Resolution of Structural Tensions.
\newblock Preprint, 2025.

\bibitem{masser-ogrady}
D.~Masser and A.~O'Grady.
\newblock The ABC Conjecture: A Survey.
\newblock {\em Bull. LMS}, 2020.

\bibitem{mochizuki-iut}
S.~Mochizuki.
\newblock Inter-universal Teichmüller theory I–IV.
\newblock Preprints, 2012–2017.

\end{thebibliography}

\end{document}
